\scanpage{40}\mspage{2}

(utilisant une orientation de $I$) qu'il existe sur chaque $V_i$ une structure de bitorseur \inserted{\unsure{ens|sur|sous}} sous une paire de groupes $G_{\underset{\substack{\Vert\\\rho^{-1}(i)}}{j}},G_{\underset{\substack{\Vert\\\rho(i)}}{k}}$ (les $G_j$ étant définis \deleted{\unsure{p}} à isom unique près)

\diagram{\begin{tikzcd}[ampersand replacement=\&,column sep=small,row sep=small]
 \& k\;\;V_k\arrow[dl,dash,"G_j"']\arrow[dr,dash,"G_i"]\&\\
 V_i\;\;i\arrow[rr,dash,"G_k"']\&\&j\;\;V_j
\end{tikzcd}}

et un isom (choisissant un $i\in I$ pour la commodité de l'écriture

\[
V_i\mathbin{\underset{G_k\deleted{\unsure{i}}}{\wedge}}V_j
 \mathbin{\underset{G_i}{\wedge}}V_k\simeq\mathbb{1}_{G_i}
 \qquad\text{isom de bitorseurs}
\]

i.e. un élément central \inserted{$\varepsilon_i$} de $V_i\mathbin{\underset{G_k}{\wedge}}V_j\mathbin{\underset{G_i}{\wedge}}V_k$ de telle façon que $V$ s'identifie à l'ens des $(x_i,x_j,x_k)$ tels que

\[
 x_i\wedge x_j\wedge x_k=\varepsilon
\]

Si les $G_i$ sont \deleted{commutatifs, ils sont} \inserted{\unsure{can}$^{\mathrm t}$. ext. \unsure{isop.}}, donc \unsure{si|où} l'un est comm, les autres le sont et on a un syst. transitif d'isomorphismes ; On est ramené à la situation précédente, dès qu'on se donne une ``origine'' $o\in V\ldots$

% archive sheet number: 39
